\section{CV 的十年学习史：从模型 scaling 到 NLP 范式迁移}

上一章建立问题，本章回到历史。张祥雨把自己的早期学术主线放在 2012 年后的深度学习 scaling 上：ImageNet 提供大数据，CUDA/GPU 提供算力，AlexNet 证明模型扩大可以显著提升效果，随后 ResNet 等工作进一步打开深层视觉网络。这个阶段的核心是 model scaling：让视觉模型更深、更大、更可训练。

后来 NLP 的 Transformer、BERT、GPT 成为新范式，CV 社区开始把这些方法迁移过来。ViT 把 Transformer 放进视觉；iGPT 尝试像语言一样自回归建模图像；MIM 和对比学习试图通过遮蔽、增强和不变性学习视觉表征。但张祥雨的判断是，许多方法只是把 NLP 的外形搬到了 CV，未必解决静态图像的根本差异。

\begin{knowledgebox}{术语消化：CV 向 NLP 学了什么}
\noindent\begin{tabularx}{\textwidth}{p{0.18\textwidth}p{0.35\textwidth}X}
\toprule
术语 & 机制 & 本期中的意义 \\
\midrule
ImageNet & 大规模标注图像数据集 & 让视觉模型第一次有稳定 scaling 数据基础。 \\
ResNet & 残差连接支持深层网络训练 & 证明视觉 model scaling 可行。 \\
ViT & 把图像切成 patch，用 Transformer 建模 & 架构迁移成功，但不等于目标函数迁移成功。 \\
iGPT/MIM & 自回归或遮蔽式图像建模 & 试图复刻语言建模，但受图像数据属性限制。 \\
Contrastive Learning & 通过增强视图学习不变性 & 很依赖人为设计的 augmentation。 \\
\bottomrule
\end{tabularx}
\end{knowledgebox}

\begin{warningbox}{课堂提示：不要把“用了 Transformer”当作 GPT 时刻}
张祥雨的提醒是，架构只是大方向。ViT、MIM、iGPT 让 CV 学到 NLP 的工具，但静态图像是不是能用同一套目标函数学出生成、理解和对齐，才是真问题。
\end{warningbox}

\subsection{静态图像为什么不是自然语言}

本节进入第一条关键论证。自然语言语料本身是人类生成的，它已经把人类概念、意图、语义和对齐关系写进 token 序列。GPT 在语言上做生成建模时，生成能力、理解能力和人类对齐能在同一目标里相互增强。静态图像则不同：图像是自然世界和成像系统产生的，它不因为人类如何理解而改变。

这就导致三个能力在图像上天然割裂。生成图像可以学像素分布；理解图像需要引入人类语义；人类对齐需要标签、caption、问答、偏好或任务定义。即使模型能生成逼真图像，也不代表它理解场景；即使模型能理解图像，也不代表它能控制生成过程。

\begin{figure}[H]
\centering
\includegraphics[width=0.92\textwidth]{static-image-three-way-split.png}
\caption{静态图像的三重割裂：生成、理解和人类对齐在静态图像中天然分离。自制概念图，依据 00:18:22--00:24:23 对谈内容整理。}
\end{figure}

\begin{knowledgebox}{读图：图像在那里，理解是人加进去的}
图中心是 Image，周围是生成、理解、对齐、自然和任务。读这张图时要抓住一句话：图像不天然等于人类语义。杯子、桌子、动作和空间关系需要人类任务框架注入，不能只从像素分布里自动得到 GPT 式语义统一。
\end{knowledgebox}

\subsection{本章小结}

CV 的十年进展说明，数据、算力、架构和 scaling 都很重要；但它也说明，视觉不是语言的简单替身。静态图像把生成、理解和对齐拆开，多模态必须先处理这个数据形态差异。

